\section{Kapsejlads}
Siden 1991 har en af de største begivenheder på Aarhus Universitet været den famøse \textit{Kapsejlads}. En begivenhed, hvor de forskellige festforeninger laver en ølstafet hen over Universitetssøen. Over tiden har det udviklet sig fra en simpel konkurence til en éndags-festival med over 30.000 ``tilskuere". \TKET har været med til arrangementet på mange måder over årene. 

\subsection{2016-2019}
I trit med de forrige år har \TKET deltaget ved kapsejladsen og hvert eneste år har vores deltagelse fulgt samme forløb: Vi vinder kapsejladsen overlegent. Det kan godt være at et af de andre holde altid ender med at få det gyldne bækken, men dette kommer sig af at de har misforstået hvad hele konkurrencen går ud på, nemlig at have det sjovest. I stræben efter at opnå dette ædle mål, går \TKET altid den ekstra mil, hvilket for det meste resultere i en diskvalifikation.

Gennem årene divergerede Umbis formål med kapsejladsen og den rigtige sejersbetingelse, hvilket desvæære havde indflydelse på \TKETs tilstedeværelse i unisøen. Da de andre atleter efterstræbede sig at komme så strømlinet igennem vandet som muligt, vagte det nemlig utilfredshed at \TKETs alternative søfartøj skabte højt farevand i søen. BEST besluttede sig for at give de andre hold en chance og lade \TKET indgå i nye roller. Fra 2017 og frem førte dette til at \TKET prøvede krafter med at være dommere og at have vores eget heat nummer $\pi$. I 2020 var \TKET det eneste hold der sejlede i unisøen, da de andre hold udeblev af urelaterede årsager.

De urelaterede årsager viste sig at være ret påtrængende og ret hurtigt blev de relaterede til alting. Corona lukkede for alle større sociale arrangementer og da kapsejladsen er et socialt arrangement i størrelsesklassen stor, blev dette også aflyst i 2021. \TKET holdt noget af kampgejsten ved lige ved at udfordre Humbug til tvekamp hvor BEST medlemmer dystede mod Humbug medlemmer i et udvalg af ukendte udfordringer.

\subsection{2022-2025}
Efter den underlige tid, kom Kapsejladsen stærkt tilbage i 2022. Dette år var vi ikke valgt til dommer eller til at sejle i heat $pi$. Så \TKETs tilstedeværelse var i \TKETs Flåde på Unisøen. Dog var det ikke en kedelig dag, da der skulle hentes øl fra breden. Efter en streng af uheldige beslutninger, endte det med, at transportflåden væltede med seks af de otte kasser drikkevarer, der var med på overfarten, og de faldt direkte i søen\footnote{Indtil 2026 hvor søen blev drænet og tømt for slam, gik en lille gruppe hvert år ned til søen for at tjekke, om øllene på bunden stadig var gode}. Selvom man kun fik reddet noget af det og derfor manglede lidt drikkevarer, var det et pænt comeback.

I 2023 lavede man en aftale med Umbilicus, med at holde et stort introshow til det daværende Kapsejlads. Lidt som tidligere år, lavede man et show, med sejlads, bunde øl, og god stemning. For at gøre dette sjovt valgte man, at der ikke skulle bygges en tømmerflåde til søen, men i stedet lavede man en fælles \TK-camp til alle de studerende under foreningen. Begge dele var en stor succes, og campen har været en ting lige siden.

I 2024 droppede man det store Introshow, og byggede igen en flåde til dagen. Dog havde man glemt at koordinere en livredder til flåden. Så flåden måtte lukkes midt på dagen. 

Det viste sig at være et større problem for i 2025. Fordi påsken lå i ugen før Kapsejlads, skulle man konstruerer Flåden i påske-weekenden. Da man så manglede en del til transportflåden, kunne man ikke købe den, da alle butikker var lukket. Derfor lavede man en nødløsning, og håbede på det bedste. Da dagen kom, og man satte flåden og transportflåden til søs, gik der ikke lang tid før man fandt ud af at transportflåden var ret usikker. Og efter et par våde sko og et par stykker var faldt i, faldt transportflåden også fra hinanden. Til sidst måtte man erkende, at det var en tabt sag, og hele flåden blev sejlet til lands. Siden da (ikke kun på grund af situationen) har man valgt ikke at have en tømmerflåde længere. Både halvdelen af bygge-træet er smidt ud, og de blå tønder er solgt. Dog har man i 2026 stadig metal-stængerne til skelettet. 

\subsection{2026}
Efter meget snak og lyst af \TKETs studerende, valgte man, at til Kapsejlads 2026 at deltage på ligefod (og på samme måde) som alle andre festforeninger på AU. Da BEST ikke har al tid i verden, lavede man en større end normal \textit{KapselKoor} gruppe, der ville stå for alt det tilhørende arbejde. Det betyder, at \TKET skulle sætte et sejlerhold, købe en båd og der blev endda også skrevet en sang\footnote{Sangen hedder ``Vi tilbage" og kan findes hvor sange høres. Mange af de andre foreningerne var meget begejstret for sangen, og var en af de mest afspillede til oprydning om lørdagen.}. Da man ikke længere havde en tømmerflåde, valgte man at navngive båden for ``Flåden", så man stadig havde en flåde på vandet. Til selve konkurrencen endte vi på en flot 10. plads!
